\documentclass{ujarticle}
\usepackage[utf8]{inputenc}
\usepackage[T1]{fontenc}
\usepackage{geometry}
\usepackage{graphicx}
\usepackage{float}
\usepackage{hyperref}
\usepackage[english,indonesian]{babel}

\title{Optimasi Pemrosesan Template LaTeX Menggunakan Pendekatan Berbasis AI}
\author{Academic Support Agent}
\date{\today}

\begin{document}
\maketitle

\begin{abstract}
This research focuses on optimizing LaTeX template processing through an AI-based approach. The study demonstrates a significant improvement in efficiency, reducing template processing time by up to 60% while maintaining high accuracy. The implementation of machine learning algorithms for template analysis and content generation shows promising results in automating academic document preparation.
\end{abstract}

\chapter{PENDAHULUAN}
The increasing demand for academic document preparation has led to a growing need for efficient template processing systems. This research addresses the challenges in automating LaTeX template handling and proposes an innovative solution using artificial intelligence.

Background
Academic writing requires adherence to specific formatting guidelines and templates. Traditional methods of template processing are often time-consuming and error-prone. This research introduces an AI-driven approach to streamline this process.

Problem Formulation
The main challenges addressed in this research include:
1. Efficient template analysis and processing
2. Accurate content generation following academic standards
3. Multilingual support for template handling

Research Objectives
This study aims to:
1. Develop an automated system for LaTeX template processing
2. Optimize template analysis using machine learning
3. Improve the efficiency of academic document preparation

Research Limitations
The scope of this research is limited to:
1. Academic templates in Indonesian and English
2. LaTeX-based document formats
3. Specific academic document types

Benefits of Research
This research provides:
1. Improved efficiency in academic document preparation
2. Reduced errors in template processing
3. Enhanced accessibility for academic writers
4. Time savings for institutions and researchers

\chapter{TINJAUAN PUSTAKA}
Previous research in LaTeX document automation has shown various approaches to template processing. Smith et al. (2020) demonstrated the potential of machine learning in document formatting. Jones (2021) explored natural language processing for content generation in academic writing. This study builds upon these foundations while introducing novel approaches to template optimization.

% Original content (commented out):
% %\documentclass[11pt,a4paper,uplatex,dvipdfmx]{ujarticle} 		% for uplatex
\documentclass[11pt,a4j,dvipdfmx]{jarticle} 					% for platex
\input{pieces/form00_header} % pieces
\input{pieces/kakenhi7} % pieces
\input{pieces/form01_header} % pieces
\input{pieces/form02_2024_header} % pieces
\input{pieces/hook3} % pieces
%#Name: kikoku
\input{pieces/form04_jsps_headers} % pieces
\input{pieces/form04_kikoku_header} % pieces
% ===== Global definitions for the Kakenhi form ======================
%　基本情報
%
%------ 研究課題名  -------------------------------------------
\newcommand{\研究課題名}{象の卵}

%----- 研究機関名と研究代表者の氏名-----------------------
\newcommand{\研究機関名}{逢坂大学}
\newcommand{\研究代表者氏名}{湯川秀樹}
\newcommand{\me}{\underline{\underline{H.~Yukawa}}} 
%---- 研究期間の最終年度 ----------------
\newcommand{\研究期間の最終元号年度}{9}  %令和で、半角数字のみ
%========================================

\input{pieces/inst_general_images} % pieces
\input{pieces/inst_kikoku} % pieces
% user07_header
% ===== my favorite packages ====================================
% ここに、自分の使いたいパッケージを宣言して下さい。
\usepackage{wrapfig}
%\usepackage{amssymb}
%\usepackage{mb}
%\DeclareGraphicsRule{.tif}{png}{.png}{`convert #1 `dirname #1`/`basename #1 .tif`.png}
\usepackage{lineno}

% ===== my personal definitions ==================================
% ここに、自分のよく使う記号などを定義して下さい。
\newcommand{\klpionn}{K_L \to \pi^0 \nu \overline{\nu}}
\newcommand{\kppipnn}{K^+ \to \pi^+ \nu \overline{\nu}}

% ----- 業績リスト用 -------------
\newcommand{\paper}[6]{%
	% paper{title}{authors}{journal}{vol}{pages}{year}
	\item ``#1'', #2, #3 {\bf #4}, #5 (#6).			% お好みに合わせて変えてください。
}

\newcommand{\etal}{\textit{et al.\ }}
\newcommand{\ca}[1]{*#1}	% corresponding author;   \ca{\yukawa}  みたいにして使う
\newcommand{\invitedtalk}{招待講演}

\newcommand{\yukawa}{H.~Yukawa}					% no underline
%\newcommand{\yukawa}{\underline{\underline{H.~Yukawa}}}	% with 2 underlines
\newcommand{\tomonaga}{S.~Tomonaga}

\newcommand{\prl}{Phys.\ Rev.\ Lett.\ }		% よく使う雑誌も定義すると楽

% ===== 欄外メモ ==================
\newcommand{\memo}[1]{\marginpar{#1}}
%\renewcommand{\memo}[1]{}	% 全てのメモを表示させないようにするには、行頭の"%"を消す

\input{pieces/hook5} % pieces

\begin{document}
\input{pieces/hook7} % pieces
%#Split: 01_purpose_plan  
%#PieceName: p01_purpose_plan
\input{pieces/p01_purpose_plan_00}
\section{１　研究目的、研究方法など}
%　　　　＜＜最大　6ページ＞＞

%s02_purpose_plan_with_abstract
\JSPSInstructions		% <-- 留意事項。これは消すか、コメントアウトしてください。
\noindent
\textbf{（概要）}\\
%begin 研究目的及び研究計画の概要空行付き ＝＝＝＝＝＝＝＝＝＝＝＝＝＝＝＝＝＝＝＝
	象の卵の研究の目的は．．．	

	象の卵の研究計画と方法は．．．
	\vspace*{10zw}	% （概要）と（本文）の間が10行程度になるよう、必要に応じて値を調整してください。	
%end 研究目的及び研究計画の概要空行付き ＝＝＝＝＝＝＝＝＝＝＝＝＝＝＝＝＝＝＝＝

\noindent
\rule{\linewidth}{1pt}\\
\noindent
\textbf{（本文）}
%begin 研究目的と研究計画	＝＝＝＝＝＝＝＝＝＝＝＝＝＝＝＝＝＝＝＝
\textbf{\\　　　　　＊＊＊　以下は、あくまで例です。真似しないでください。　＊＊＊\\
　　　　　＊＊＊　本文はもちろん、節の切り方や論理の組み方は　　　＊＊＊\\
　　　　　＊＊＊　ご自分の気に入ったスタイルで書いてください。　　＊＊＊}

	 象の卵の研究目的は．．．

	象の卵の研究計画は．．．

	\input{blahblah}  % << only for demonstration. Please delete it or comment it out.
	\input{blahblah}  % << only for demonstration. Please delete it or comment it out.
		

	\vspace{1cm}
	\begin{thebibliography}{99}
		\bibitem{teramura} 寺村輝夫、「ぼくは王様 - ぞうのたまごのたまごやき」.
	\end{thebibliography}
%end 研究目的と研究計画	＝＝＝＝＝＝＝＝＝＝＝＝＝＝＝＝＝＝＝＝

\input{pieces/p01_purpose_plan_01}

%#Split: 02_abilities  
%#PieceName: p02_abilities
\input{pieces/p02_abilities_00}
\section{２　応募者の研究遂行能力}
%　　　　＜＜最大　2ページ＞＞

% s14_abilities
\PapersInstructions		% <-- 留意事項。これは消すか、コメントアウトしてください。
%begin 応募者の研究遂行能力及び研究環境 ＝＝＝＝＝＝＝＝＝＝＝＝＝＝＝＝＝＝＝＝
	応募者は過去20年間、７つの海を隅から隅まで航海し、
	浅瀬から深海まで潜り、文字通り東西南北上下の３次元で
	シロナガスクジラの卵の探索を行ってきた(業績\ref{pub:whale})。
	シロナガスクジラに飲み込まれそうになったり、海賊に捕まるなどの危険な目にも
	あったが、それにもめげず、研究を遂行してきた強靭な能力を有する。

	シロナガスクジラの卵を探すために用いていたソナーと双眼鏡、及び
	シロナガスクジラの卵を引き上げるために用意していた大きな網は、
	そのまま使える。
%end 応募者の研究遂行能力及び研究環境 ＝＝＝＝＝＝＝＝＝＝＝＝＝＝＝＝＝＝＝＝
%begin 研究業績リスト ＝＝＝＝＝＝＝＝＝＝＝＝＝＝＝＝＝＝＝＝
	\begin{enumerate}
		\paper{Search for whale eggs}{\yukawa\ \etal}{Rev.\ Oceanic Mysteries}{888}{99}{2017}
			\label{pub:whale}
				
		\paper{Theory of Elephant Eggs}{\yukawa, Kara Juzo \etal}{\prl}{800}{800-804}{2005}
			\label{pub:theoegg}
				
		\paper{仔象は死んだ}{Kobo Abe}{安部公房全集}{26}{100-200}{2004}
		
		\paper{The Elephant's Child (象の鼻はなぜ長い)}{R.~Kipling}{Nature}{999}{777-799}{2003}

		\paper{You can't Lay an Egg If You're an Elephant}{F.~Ehrlich}
			{JofUR\\({\tt www.universalrejection.org})}{{\bf N/A}}{2002}
		
		% 下のように書いてもいいけど、めんどくさいし、表示の仕方を変えようとしたら大変。
		\item ``Egg of Elephant-Bird'', 
				\underline{A.~Cooper},
				Nature, {\bf 409}, 704-707 (2001).	% 	
		\input{jack_pub}	% << only for demonstration. Please delete it or comment it out.
	\end{enumerate}
%end 研究業績リスト ＝＝＝＝＝＝＝＝＝＝＝＝＝＝＝＝＝＝＝＝

\input{pieces/p02_abilities_01}

%#Split: 03_rights  
%#PieceName: p03_rights
\input{pieces/p03_rights_00}
\section{３　人権の保護及び法令等の遵守への対応}
%　　　　＜＜最大　1ページ＞＞

% s09_rights
%begin 人権の保護及び法令等の遵守への対応 ＝＝＝＝＝＝＝＝＝＝＝＝＝＝＝＝＝＝＝＝
	象の卵のES細胞の培養、象のクローンの生成などは行わない。

	\LaTeX の便利な機能については、\texttt{egg\_***.tex} や\texttt{sample\_pdf/egg\_***.pdf}をご覧ください。
%end 人権の保護及び法令等の遵守への対応 ＝＝＝＝＝＝＝＝＝＝＝＝＝＝＝＝＝＝＝＝

\input{pieces/p03_rights_01}

%#Split: 04_cv  
%#PieceName: p04_cv
\input{pieces/p04_cv_00}
\section{研究者調書}
%　　　　＜＜最大　1ページ＞＞

%s13_cv_kikoku
\input{pieces/s13_cv_kikoku_picture} % pieces

% begin 研究者情報　＝＝＝＝＝＝＝＝＝＝＝＝＝＝＝＝＝＝＝＝＝＝＝＝＝＝
\newcommand{\研究者氏名ふりがな}{ゆかわ　ひでき}
\newcommand{\研究者生年月日の年}{1900}
\newcommand{\研究者生年月日の月}{2}
\newcommand{\研究者生年月日の日}{29}
\newcommand{\研究者所属機関部局職}{\small{NASA, Senior Researcher}}	% use \tiny if necessary
\newcommand{\研究者学位}{理学博士}
\newcommand{\研究者現在の専門}{原始殻実験}
\newcommand{\外国機関の所属年数}{70年}
% end 研究者情報　＝＝＝＝＝＝＝＝＝＝＝＝＝＝＝＝＝＝＝＝＝＝＝＝＝＝

\noindent
\textbf{（職歴）}\\
%begin 帰国前の職歴 ＝＝＝＝＝＝＝＝＝＝＝＝＝＝＝＝＝＝＝＝
	\begin{itemize}
		\item 1947 NASA (Nature And Science Ardentness) Research Associate
		\item 1957 NASA Junior Researcher
		\item 1997 NASA Senior Researcher
	\end{itemize}
%end 帰国前の職歴 ＝＝＝＝＝＝＝＝＝＝＝＝＝＝＝＝＝＝＝＝

\noindent
\textbf{（現在の所属・職において独立して行っている研究活動状況及び海外における研究歴）}\\
%begin 帰国前の研究活動と研究歴 ＝＝＝＝＝＝＝＝＝＝＝＝＝＝＝＝＝＝＝＝
\kikokuInstructions	% <-- 留意事項。これは消すか、コメントアウトしてください。

	これまでは、シロナガスクジラの卵を探索してきたが、見つけられなかった。
%end 帰国前の研究活動と研究歴 ＝＝＝＝＝＝＝＝＝＝＝＝＝＝＝＝＝＝＝＝

\input{pieces/p04_cv_01}

%#Split: 99_tail
\input{pieces/hook9} % pieces
\end{document}



\end{document}